% ============================================================================
%  lessons/lesson01.tex — درسِ اوّل: مِنَ الْأَشْعَارِ الْمَنْسُوبَةِ إِلَى
%                          الإِمَامِ عَلِيٍّ (ع)
%  محتوا: متنِ کتاب + ترجمه + معجم + قواعد (حروفِ مشبهه، لا نافیهٔ جنس) +
%         خودآزمایی و تمرین‌های کتاب + پاسخ‌نامهٔ تشریحی
% ============================================================================

\Lesson
  {اَلدَّرْسُ الْأَوَّلُ}
  {مِنَ الْأَشْعَارِ الْمَنْسُوبَةِ إِلَى الإِمَامِ عَلِيٍّ (ع)}
  {قِيمَةُ كُلِّ امْرِئٍ مَا يُحْسِنُهُ.}
  {ارزشِ هر انسانی به [همان چیزی] است که آن را نیکو انجام می‌دهد.}
  {اَمِيرُ الْمُؤْمِنِينَ عَلِيٌّ (ع)}

\JGoals{شناخت و ترجمهٔ حروفِ مشبههٔ فعل (\emph{إِنَّ، أَنَّ، كَأَنَّ، لٰكِنَّ،
لَيْتَ، لَعَلَّ}) \JGoalSep معنای چهارمِ «لا» یعنی لا نافیهٔ جنس \JGoalSep
واژه‌های مُعجَمِ درس \JGoalSep ترجمهٔ آیات و احادیث با این قواعد}

% ============================================================================
%  ۱) متنِ درس
% ============================================================================
\Jsec{اَلنَّصُّ \,\textbar\, متنِ درس}

\begin{JNass}[عنوانِ درس \,\textbar\, مِنَ الْأَشْعَارِ الْمَنْسُوبَةِ إِلَى الإِمَامِ عَلِيٍّ (ع)]
شعرهای این درس به امامِ اوّلِ شیعیان، حضرتِ علی بنِ ابی‌طالب (ع) منسوب‌اند؛
حکمتِ گزینشِ آن‌ها، آشنایی با ارزش‌های والایی چون \textbf{علم}، \textbf{عفاف}
و \textbf{خودشناسی} در کلامِ امیرِ مؤمنان است.
\end{JNass}

% ------------------------------------------------------ قطعهٔ اوّل ----------
\begin{JNass}[اَلدَّاءُ وَ الدَّوَاءُ \,\textbar\, بیماری و درمان]
\begin{JPoem}
\JPoemBayt{دَواؤُكَ فيكَ وَ ما تُبْصِرُ}{وَ داؤُكَ مِنكَ وَ لا تَشْعُرُ}
\JPoemBayt{أَ تَزْعَمُ أَنَّكَ جِرْمٌ صَغير}{وَ فيكَ انْطَوَى الْعالَمُ الْأَكْبَرُ}
\end{JPoem}
\end{JNass}

\begin{JTarjome}[ترجمهٔ بیت‌ها]
در توضیحِ شعر: \textbf{دَوَاء} به معنای «بیماری» است و \textbf{دَوَاؤُك} یعنی
«بیماری‌ات». معنای سروده چنین است:

درمان‌ها (بیماری‌ها) در توست و نمی‌بینی؛ و درمانِ (داروی) تو از خودِ توست و
احساسش نمی‌کنی. آیا می‌پنداری که تو پیکری کوچکی، در حالی‌که جهانِ بزرگ‌تر در
تو به هم پیچیده (نهان) شده است؟

\noindent{\JUIFont\footnotesize\color{JAccentD}\textbf{پیامِ بیت:}}
انسان از بزرگ‌ترین نعمتِ خدادادی یعنی «خود» غافل است؛ در حالی‌که شناختِ خویشتن،
کلیدِ همهٔ شناخت‌ها و درمانِ درماندگی‌هاست.
\end{JTarjome}

% ------------------------------------------------------ قطعهٔ دوم -----------
\begin{JNass}[اَلنَّاسُ أَكْفَاءُ \,\textbar\, مردم یکسان‌اند]
\begin{JPoem}
\JPoemBayt{اَلنَّاسُ مِن جِهَةِ الْآبَاءِ أَكْفَاءُ}{أَبُوهُمُ آدَمُ وَ الْأُمُّ حَوَّاءُ}
\JPoemBayt{وَ قَدْرُ كُلِّ امْرِئٍ مَا كَانَ يُحْسِنُهُ}{وَ لِلرِّجَالِ عَلَى الْأَفْعَالِ أَسْمَاءُ}
\JPoemBayt{فَفُزْ بِعِلْمٍ وَ لا تَطْلُبْ بِهِ بَدَلاً}{فَالنَّاسُ مَوْتَى وَ أَهْلُ الْعِلْمِ أَحْيَاءُ}
\end{JPoem}

\JPNote{١}{تنوینِ کلمهٔ «آدَمٍ» و سکونِ قافیه به دلیلِ ضرورتِ شعری است.}
\JPNote{٢}{و ارزشِ هر انسانی به کاری است که آن را خوب انجام می‌دهد و انسان‌ها
با کارهایشان نام‌هایی دارند. (انسان‌ها با کارهایشان شناخته شده و با آن نامیده
می‌شوند.)}
\JPNote{٣}{پس دانش را به دست بیاور و برایش جانشینی نخواه که مردم، مُرده و
اهلِ دانش زنده‌اند.}
\end{JNass}

% ------------------------------------------------------ قطعهٔ سوم -----------
\begin{JNass}[اَلْفَخْرُ بِالْعَفَافِ \,\textbar\, افتخار به بی‌نیازی]
\begin{JPoem}
\JPoemBayt{أَيُّهَا الْفَاخِرُ جَهْلاً بِالنَّسَبِ}{إِنَّمَا النَّاسُ لِأُمٍّ وَ لِأَبِ}
\JPoemBayt{هَلْ تَرَاهُم خُلِقُوا مِن فِضَّةٍ}{أَمْ حَدِيدٍ أَمْ نُحَاسٍ أَمْ ذَهَبْ}
\JPoemBayt{بَلْ تَرَاهُم خُلِقُوا مِن طِينَةٍ}{هَلْ سِوَى لَحْمٍ وَ عَظْمٍ وَ عَصَبْ}
\JPoemBayt{إِنَّمَا الْفَخْرُ لِعَقْلٍ ثَابِتٍ}{وَ حَيَاءٍ وَ عَفَافٍ وَ أَدَبْ}
\end{JPoem}
\end{JNass}

% ============================================================================
%  ۲) ترجمهٔ روان
% ============================================================================
\Jsec{اَلتَّرْجَمَةُ \,\textbar\, ترجمهٔ روانِ درس}

\begin{JTarjome}[ترجمهٔ اَلنَّاسُ أَكْفَاءُ]
مردم از نظرِ پدران یکسان‌اند؛ پدرِ همه آنان آدم و مادرشان حوا است.

و ارزشِ هر انسانی به آن چیزی است که آن را نیکو انجام می‌دهد و مردان به
پاسِ کارهایشان، نام‌هایی دارند.

پس با علم، پیروز و بهره‌مند شو و از آن بهرهٔ مادی مجوی؛ که مردم مُرده‌اند و
اهلِ علم زنده‌اند.
\end{JTarjome}

\begin{JTarjome}[ترجمهٔ اَلْفَخْرُ بِالْعَفَافِ]
ای آن‌که به نادانی به تبار و نژاد (پدر و مادر) خویش فخر می‌فروشی، همهٔ مردم
از یک مادر و یک پدرند.

آیا می‌بینی آنان از نقره آفریده شده‌اند یا آهن، یا مس، یا طلا؟

بلکه می‌بینی که آنان از گِل (سرشتی) آفریده شده‌اند؛ آیا جز گوشت و استخوان و
پی (عصب) چیزی دیگرند؟

افتخارِ حقیقی تنها برای خردی استوار و حیاء و پرهیزکاری و ادب است.
\end{JTarjome}

\begin{JGood}[چرا این شعرها به امیرِ مؤمنان (ع) منسوب‌اند؟]
دورةٔ شعرِ عربیِ متقدّم، سروده‌های بسیاری با سیمای «حِکَم» (حکمت‌ها و
پندنامه‌ها) به امام علی (ع) نسبت داده است. محتوای این سروده‌ها — دعوت به
خودشناسی، برابریِ انسان‌ها، ارجِ علم و دوری‌جویی از تکبّرِ نژادی — کاملاً
هماهنگ با آموزه‌های کلامِ آن حضرت در نهج‌البلاغه است.
\end{JGood}

% ============================================================================
%  ۳) واژه‌نامهٔ درس
% ============================================================================
\Jsec{اَلْمُعْجَمُ \,\textbar\, واژه‌نامهٔ درس}

\begin{JMTable}[اَلْمُعْجَمُ \,\textbar\, واژه‌نامهٔ درسِ اوّل]
\JW{اِنْطَوَى}{به هم پیچیده شد، نهان شد}{مضارع: يَنْطَوِي}
\JW{أَبْصَرَ}{نگاه کرد}{}
\JW{اَلْبَدَل}{جانشین، بهرهٔ مادی}{جمع: اَلْأَبْدَال}
\JW{اَلْجِرْم}{پیکر، تن}{جمع: اَلْأَجْرَام}
\JW{اَلْحَدِيد}{آهن}{}
\JW{اَلدَّاء}{بیماری}{\textbf{=} اَلْمَرَض؛ \textbf{≠} اَلشِّفَاء، اَلصِّحَّة}
\JW{زَعَمَ}{گمان کرد}{}
\JW{سِوَى}{جُز، غیر از}{}
\JW{اَلطِّين ، اَلطِّينَة}{گِل، سرشت}{}
\JW{اَلْعَصَب}{پی، عصب}{جمع: اَلْأَعْصَاب}
\JW{اَلْعَظْم}{استخوان}{جمع: اَلْعِظَام}
\JW{اَللَّحْم}{گوشت}{جمع: اَللُّحُوم}
\JW{اَلنُّحَاس}{مس}{}
\end{JMTable}

\begin{JNokte}[تفاوتِ نزدیک‌به‌هم]
\textbf{اَلدَّاء} (بیماری) با \textbf{اَلدُّعَاء} (دعا) اشتباه نشود؛ در ترجمه،
«دَوَاؤُكَ» یعنی \emph{بیماری‌ات}، نه «داروی تو»! همین نکتهٔ ظریف، موضوعِ بیت
نخست را شکل می‌دهد: «بیماری‌ات در توست و نمی‌بینی».
\end{JNokte}

% ============================================================================
%  ۴) قواعدِ درس
% ============================================================================
\Jsec{إِعْلَمُوا \,\textbar\, آموزشِ قواعد}

\Jsub{معانیِ حروفِ مشبههٔ فعل و «لا»ی نافیهٔ جنس}

% ------------------------------------------------ ۴.۱ حروف مشبهه بالفعل -----
\Jsubsec{اَلْحُرُوفُ الْمُشَبَّهَةُ بِالْفِعْلِ}

با معانیِ دقیقِ حروفِ پرکاربردِ «إِنَّ ، أَنَّ ، كَأَنَّ ، لٰكِنَّ ، لَيْتَ ،
لَعَلَّ» آشنا شوید:

\JHadithCard{إِنَّ}{قطعاً، همانا، به‌درستی‌که، بی‌گمان}
{(... إِنَّ اللّٰهَ لا يُضيعُ أَجْرَ الْمُحْسِنينَ)}{اَلتَّوْبَة: ١٢٠}
{بی‌گمان خدا پاداشِ نیکوکاران را تباه نمی‌کند.}
{جملهٔ پس از خود را تأکید می‌کند.}

\JHadithCard{أَنَّ}{که}
{(... قالَ أَعْلَمُ أَنَّ اللّٰهَ عَلَىٰ كُلِّ شَيْءٍ قَدِيرٌ)}{اَلْبَقَرَة: ٢٥٩}
{گفت: می‌دانم \emph{که} خدا بر هر چیزی تواناست.}
{دو جمله را به هم پیوند می‌دهد.}

\begin{JNokte}[اَلْمُرَكَّبُ لَعَنْدَنَا \,\textbar\, «لِأَنَّ»]
ترکیبِ «\textbf{لِأَنَّ}» به معنای «\textbf{زیرا، برای این‌که}» است و در پاسخِ
پرسش‌های «چرا» می‌آید:

\noindent\begin{tabular}{@{}p{0.47\linewidth}@{\hspace{1em}}p{0.47\linewidth}@{}}
\Ar \textbf{پُرسش:} لِمَاذَا مَا سَافَرْتَ بِالطَّائِرَةِ؟ &
\Ar \textbf{پاسخ:} لِأَنَّ بِطَاقَةَ الطَّائِرَةِ غَالِيَةٌ. \\[2pt]
{\small چرا با هواپیما سفر نکردی؟} & {\small زیرا بلیتِ هواپیما گران است.} \\
\end{tabular}
\end{JNokte}

\JHadithCard{كَأَنَّ}{گویی، مانندِ}
{(كَأَنَّهُنَّ الْيَاقُوتُ وَ الْمَرْجَانُ)}{اَلرَّحْمٰن: ٥٨}
{آنان گویی یاقوت و مرجان‌اند.}
{برای تشبیه به کار می‌رود.}

\begin{JMesal}
\Ar كَأَنَّ إِرْضاءَ جَمِيعِ النَّاسِ غَايَةٌ لا تُدْرَكُ.

گویی خشنود ساختنِ همهٔ مردم، هدفی است که به دست نمی‌آید.
\end{JMesal}

\JHadithCard{لٰكِنَّ}{ولی، اما}
{(... إِنَّ اللّٰهَ لَذُو فَضْلٍ عَلَى النَّاسِ وَلٰكِنَّ أَكْثَرَ النَّاسِ لا
يَشْكُرُونَ)}{اَلْبَقَرَة: ٢٤٣}
{بی‌گمان خدا صاحبِ فضل و بخشش بر مردم است، ولی بیشترِ مردم سپاسگزاری نمی‌کنند.}
{پیامِ جملهٔ پیشین را کامل می‌کند و ابهام را برمی‌دارد.}

\JHadithCard{لَيْتَ}{کاش}
{(وَ يَقُولُ الْكَافِرُ يَا لَيْتَنِي كُنْتُ تُرَابًا)}{اَلنَّبَأ: ٤٠}
{و کافر می‌گوید: ای کاش من خاک بودم!}
{بیانگرِ آرزوست؛ به صورتِ «يَا لَيْتَ» نیز می‌آید.}

\JHadithCard{لَعَلَّ}{شاید، امید است}
{(إِنَّا جَعَلْنَاهُ قُرْآنًا عَرَبِيًّا لَعَلَّكُم تَعْقِلُونَ)}{اَلزُّخْرُف: ٣}
{بی‌گمان ما آن را قرآنی عربی قرار دادیم؛ امید است که بیندیشید.}
{نشانِ امید و انتظار است.}

% ------------------------------------------------ توضیح اعرابی --------------
\Jsubsec{جایگاهِ اعرابیِ اسم و خبرِ این حروف}

حروفِ مشبههٔ فعل بر سرِ \textbf{جملهٔ اسمیّه} (مبتدا و خبر) می‌آیند و هنگامی‌که
بر اسمِ ظاهر وارد می‌شوند، آن را به‌عنوانِ «اسمِ خود» \textbf{منصوب} می‌کنند؛
ولی اعرابِ خبر را تغییر نمی‌دهند و خبر، مرفوع می‌ماند:

\begin{JMesal}
\noindent
\begin{tabular}{@{}p{0.49\linewidth}@{\hspace{0.02\linewidth}}p{0.49\linewidth}@{}}
\multicolumn{1}{@{}c}{\JUIFont\bfseries\small\color{JSlate}پیش از ورودِ حرف} &
\multicolumn{1}{c@{}}{\JUIFont\bfseries\small\color{JSlate}پس از ورودِ حرف}\\[4pt]
\Ar مَهْدِيٌّ فَائِزٌ فِي مُسَابَقَةِ كُرَةِ الْمِنْضَدَةِ. &
\Ar لَيْتَ مَهْدِيًّا فَائِزٌ فِي مُسَابَقَةِ كُرَةِ الْمِنْضَدَةِ! \\[2pt]
{\small مَهديٌّ: مبتدا، مرفوع} & {\small مَهديًّا: اسمِ لَيْتَ، منصوب}\\
{\small فائزٌ: خبر، مرفوع} & {\small فائزٌ: خبرِ لَيْتَ، مرفوع}\\
\end{tabular}
\end{JMesal}

\begin{JDam}[دامِ آزمون]
وقتی خودِ حروفِ مشبهه به «مبتدا و خبر» می‌چسبند، علامت‌ها عوض می‌شوند:
«إِنَّ + الْعِلْمَ نَافِعٌ»؛ یعنی اسم منصوب (الْعِلْمَ) و خبر مرفوع (نَافِعٌ).
گزینه‌هایی که «الْعِلْمُ نَافِعٌ» یا «الْعِلْمَ نَافِعًا» می‌آورند، دام‌اند!
همچنین «تَذَكُّرٌ مُهِمٌّ»: در ترجمهٔ صحیح باید به سیاقِ عبارت (بافتِ سخن)
توجّه کنی؛ ترجمه، هنر و علم است و باید ذوق و سلیقهٔ مترجم همراهِ توانمندی‌های
زبانی، مانندِ شناختِ ویژگی‌های زبانِ مبدأ و مقصد باشد.
\end{JDam}

% ------------------------------------------------ ۴.۲ لا نافیه جنس ---------
\Jsubsec{لا النَّافِيَةُ لِلْجِنْسِ (لای نفی جنس)}

تاکنون با سه معنای حرفِ «لا» آشنا شده‌ای:
\begin{enumerate}
  \item «لا» به معنای «نه» در پاسخ به «هَلْ» و «أَ»؛ مانندِ:
        أَ أَنْتَ مِنْ بُجْنُورد؟ لا، أَنَا مِنْ بِيرْجَنْد.
  \item لایِ نفیِ مضارع؛ مانندِ: لا يَذْهَبُ: نمی‌رود.
  \item لایِ نهی؛ مانندِ: لا تَذْهَبْ: نرو؛ و به معنای «نباید» بر سرِ فعلِ
        مضارعِ اوّل و سومِ شخص؛ مانندِ: لا يَذْهَبْ: نباید برود.
\end{enumerate}

در این درس با معنای چهارمِ «لا» آشنا می‌شویم:

\begin{JQaede}[JCompact]{قاعدهٔ لا نافیهٔ جنس}
معنای چهارمِ «لا»، «\textbf{هیچ... نیست}» است و به آن «\textbf{لای نفی جنس}»
می‌گویند؛ این «لا» بر سرِ \textbf{اسم} وارد می‌شود:

\begin{itemize}
  \item \textbf{اسمِ لا:} معمولاً \textbf{بدونِ «ال»} و دارایِ \textbf{فتحه} است.
  \item \textbf{خبرِ لا:} اگر اسمِ ظاهر (آشکار) باشد، به همان شکل
        \textbf{مرفوع} می‌ماند؛ و اغلب به صورتِ \textbf{جارّ و مجرور} می‌آید.
\end{itemize}

\Ar لا عِلْمَ لَنَا إِلَّا مَا عَلَّمْتَنَا ... (اَلْبَقَرَة: ٣٢)

جز آنچه به ما آموخته‌ای، هیچ دانشی نداریم. (عِلْمَ: اسمِ لا، بدونِ ال و دارای
فتحه)
\end{JQaede}

\begin{JHadith}
\Ar لا مُصِيبَةَ أَعْظَمُ مِنَ الْجَهْلِ. اَلْإِمَامُ الصَّادِقُ (ع)

هیچ بلايی بزرگ‌تر از نادانی نیست.
\end{JHadith}

سه پرسشِ کتاب را با هم مرور کنیم:
\begin{enumerate}
  \item اسمِ بعد از این لا چه حرکتی دارد؟ \textcolor{JGreen}{\textbf{فتحه دارد.}}
  \item آیا اسمِ پس از لا دارایِ «ال» است؟ \textcolor{JGreen}{\textbf{خیر؛
        معمولاً بدونِ ال است.}}
  \item پس می‌توان گفت: اسمِ پس از لایِ نفیِ جنس «بدونِ ال و مفتوح» است؟
        \textcolor{JGreen}{\textbf{آری.}}
\end{enumerate}

\begin{JMesal}
خبرِ لایِ نفیِ جنس در صورتی‌که اسمِ ظاهر (آشکار) باشد، به همان شکل مرفوع
می‌ماند؛ مثال:

\noindent
\begin{tabular}{@{}p{0.48\linewidth}@{\hspace{0.02\linewidth}}p{0.48\linewidth}@{}}
\Ar اَلْعَفْوُ عِنْدَ الْقُدْرَةِ حَسَنٌ. &
\Ar لا شَيْءَ أَحْسَنُ مِنَ الْعَفْوِ عِنْدَ الْقُدْرَةِ.\\[2pt]
{\small} & {\small شَيْءَ: اسمِ لا، بدونِ «ال» و دارایِ فتحه} \\
{\small} & {\small أَحْسَنُ: خبرِ لا، مرفوع} \\
\end{tabular}
\end{JMesal}

\begin{JGood}
خبرِ لای نفی جنس اغلب به صورتِ «جارّ و مجرور» می‌آید؛ مثال:

\Ar لا رَجُلَ فِي الْحَفْلَةِ.\quad لا مَاءَ فِي الْبَيْتِ.\quad
لا خَائِنَ فِي فَرِيقِنَا.\quad لا شَكَّ فِيهِ.
\end{JGood}

\begin{JNokte}[حذفِ خبر]
گاهی خبرِ لای نفی جنس حذف می‌شود؛ مانندِ: لا إِلٰهَ [مَوْجُودٌ] إِلَّا اللّٰهُ
و لا شَكَّ و لا بَأْسَ.
\end{JNokte}

% ============================================================================
%  ۵) خودآزمایی (بدونِ پاسخ)
% ============================================================================
\Jsec{اِخْتَبِرْ نَفْسَكَ \,\textbar\, خودآزماییِ درس}

\begin{JTest}[خودآزماییِ ١ \,\textbar\, ترجمهٔ آیات]
تَرْجِمْ هَاتَيْنِ الْآيَتَيْنِ الْكَرِيمَتَيْنِ (این دو آیهٔ کریمه را ترجمه کن):

\begin{enumerate}
  \item ﴿فَهٰذَا يَوْمُ الْبَعْثِ\FNmark{١} وَلٰكِنَّكُم كُنْتُمْ لا تَعْلَمُونَ﴾
        اَلرُّوم: ٥٦
        \FNfoot{١}{اَلْبَعْث: رستاخیز}
  \item ﴿إِنَّ اللّٰهَ يُحِبُّ الَّذِينَ يُقَاتِلُونَ فِي سَبِيلِهِ صَفًّا\FNmark{٢}
        كَأَنَّهُم بُنْيَانٌ مَرْصُوصٌ\FNmark{٣}﴾ اَلصَّف: ٤
        \FNfoot{٢}{صَفًّا: صف‌درصف} \FNfoot{٣}{اَلْبُنْيَانُ الْمَرْصُوص:
        ساختمانِ استوار}
\end{enumerate}
\end{JTest}

\begin{JTest}[خودآزماییِ ٢ \,\textbar\, اسم و خبرِ حروفِ مشبهه]
عَيِّنِ اسْمَ الْحُرُوفِ الْمُشَبَّهَةِ وَ خَبَرَهَا؛ ثُمَّ اذْكُرْ إِعْرَابَهُمَا
(اسمِ حروفِ مشبهه و خبرشان را مشخص کن؛ سپس اعرابِ آن دو را ذکر کن):

\begin{enumerate}
  \item لَيْتَ فَصْلَ الرَّبِيعِ طَوِيلٌ فِي بَلَدِنَا! لِأَنَّ الرَّبِيعَ
        قَصِيرٌ هُنَا.
  \item كَأَنَّ الْمُشْتَرِيَ مُتَرَدِّدٌ فِي شِرَاءِ الْبِضَاعَةِ؛ وَلٰكِنَّ
        الْبَائِعَ عَازِمٌ عَلَىٰ بَيْعِهَا.
  \item اِبْحَثْ عَنْ مَعْنَى «اَلْعَصَّارَةِ\FNmark{١}» فِي الْمُعْجَمِ؛ لَعَلَّ
        الْكَلِمَةَ مَكْتُوبَةٌ فِيهِ!
        \FNfoot{١}{اَلْعَصَّارَة: آب‌میوه‌گیری}
\end{enumerate}
\end{JTest}

\begin{JTest}[خودآزماییِ ٣ \,\textbar\, ترجمهٔ احادیث]
تَرْجِمْ هٰذِهِ الْأَحَادِيثَ حَسَبَ قَوَاعِدِ الدَّرْسِ، ثُمَّ عَيِّنِ اسْمَ لا
النَّافِيَةِ لِلْجِنْسِ وَ خَبَرَهَا:

\begin{enumerate}
  \item لا خَيْرَ فِي قَوْلٍ إِلَّا مَعَ الْفِعْلِ. رَسُولُ اللّٰهِ (ص)
  \item لا جِهَادَ كَجِهَادِ النَّفْسِ. أَمِيرُ الْمُؤْمِنِينَ عَلِيٌّ (ع)
  \item لا لِبَاسَ أَجْمَلُ مِنَ الْعَافِيَةِ. أَمِيرُ الْمُؤْمِنِينَ عَلِيٌّ (ع)
  \item لا فَقْرَ كَالْجَهْلِ وَ لا مِيرَاثَ كَالْأَدَبِ. أَمِيرُ الْمُؤْمِنِينَ
        عَلِيٌّ (ع)
  \item لا سُوءَ أَسْوَأُ مِنَ الْكَذِبِ. أَمِيرُ الْمُؤْمِنِينَ عَلِيٌّ (ع)
\end{enumerate}
\end{JTest}

\begin{JTest}[خودآزماییِ ٤ \,\textbar\, جایِ خالی و نوعِ «لا»]
اِمْلَأِ الْفَرَاغَ فِيمَا يَلِي\FNmark{١}، ثُمَّ عَيِّنْ نَوْعَ «لا» فِيهِ
(جای خالی را پر کن و نوعِ «لا» را مشخص کن):

\begin{enumerate}
  \item ﴿... هَلْ يَسْتَوِي الَّذِينَ يَعْلَمُونَ وَ الَّذِينَ لا يَعْلَمُونَ
        ...﴾ اَلزُّمَر: ٩\\
        آیا کسانی‌که می‌دانند و کسانی‌که \JdotsW{4cm} برابرند؟
  \item ﴿وَ لا يَحْزُنْكَ قَوْلُهُمْ إِنَّ الْعِزَّةَ لِلّٰهِ جَمِيعًا ...﴾
        يُونُس: ٦٥\\
        گفتارشان تو را \JdotsW{3.2cm}؛ زیرا ارجمندی، همه \JdotsW{2.8cm} خداست.
  \item ﴿ذٰلِكَ الْكِتَابُ لا رَيْبَ\FNmark{٢} فِيهِ هُدًى لِلْمُتَّقِينَ﴾
        اَلْبَقَرَة: ٢\\
        آن کتاب \JdotsW{4.2cm} هدایتی برای پرهیزگاران است.
        \FNfoot{٢}{اَلرَّيْب: شک}
  \item ﴿لا إِكْرَاهَ فِي الدِّينِ ...﴾ اَلْبَقَرَة: ٢٥٦\\
        \JdotsW{3.5cm} در دین \JdotsW{3.5cm}.
  \item لا تُطْعِمُوا الْمَسَاكِينَ مِمَّا لا تَأْكُلُونَ. رَسُولُ اللّٰهِ (ص)\\
        \Jdots
\end{enumerate}
\FNfoot{١}{مَا يَلِي: آنچه می‌آید}
\end{JTest}

% ============================================================================
%  ۶) تمرین‌های کتاب (بدونِ پاسخ)
% ============================================================================
\Jsec{اَلتَّمَارِينُ \,\textbar\, تمرین‌های کتاب}

\begin{JTamrin}{١}{جایِ خالی با واژه‌های مُعجَم}
ضَعْ فِي الْفَرَاغِ كَلِمَةً مُنَاسِبَةً لِلتَّوْضِيحَاتِ التَّالِيَةِ مِنْ
مُعْجَمِ الدَّرْسِ (با توجه به توضیحِ هر سطر، واژهٔ مناسب از معجمِ درس را
در جای خالی بنویس):

\begin{enumerate}
  \item اَلـ\JdotsW{3cm} عُنْصُرٌ فِلِزِّيٌّ كَالْحَدِيدِ مُوَصِّلٌ
        لِلْحَرَارَةِ وَ الْكَهْرَبَاءِ. (اَلْمُوَصِّل: رسانا)
  \item اَلـ\JdotsW{3cm} مَادَّةٌ حَمْرَاءُ مِنْ جِسْمِ الْحَيَوَانِ تُصْنَعُ
        مِنْهُ أَطْعِمَةٌ. (اَلْحَمْرَاء: اَلْأَحْمَر)
  \item اَلـ\JdotsW{3cm} خَيْطٌ أَبْيَضُ فِي الْجِسْمِ يَجْرِي فِيهِ الْحِسُّ.
        (اَلْخَيْط: نخ)
  \item اَلـ\JdotsW{3cm} قِسْمٌ قَوِيٌّ مِنَ الْجِسْمِ عَلَيْهِ لَحْمٌ.
  \item اَلـ\JdotsW{3cm} تُرَابٌ مُخْتَلِطٌ بِالْمَاءِ.
\end{enumerate}
\end{JTamrin}

\begin{JTamrin}{٢}{استخراجِ مطلوب از ابیات}
اِسْتَخْرِجْ مِمَّا يَلِي الْمَطْلُوبَ مِنْكَ (موردِ خواسته را از هر بیت
استخراج کن):

\begin{enumerate}
  \item \textbf{اَلْمُبْتَدَأَ وَ إِعْرَابَهُ:}\\
        دَواؤُكَ فِيكَ وَ مَا تُبْصِرُ / وَ دَاؤُكَ مِنْكَ وَ لا تَشْعُرُ
  \item \textbf{اِسْمَ التَّفْضِيلِ وَ مَحَلَّهُ الْإِعْرَابِيَّ:}\\
        أَ تَزْعَمُ أَنَّكَ جِرْمٌ صَغِيرٌ / وَ فِيكَ انْطَوَى الْعَالَمُ
        الْأَكْبَرُ
  \item \textbf{اَلْخَبَرَ وَ إِعْرَابَهُ:}\\
        اَلنَّاسُ مِنْ جِهَةِ الْآبَاءِ أَكْفَاءُ / أَبُوهُمُ آدَمُ وَ الْأُمُّ
        حَوَّاءُ
  \item \textbf{اَلْجَمْعَ الْمُكَسَّرَ:}\\
        وَ قَدْرُ كُلِّ امْرِئٍ مَا كَانَ يُحْسِنُهُ / وَ لِلرِّجَالِ عَلَى
        الْأَفْعَالِ أَسْمَاءُ
  \item \textbf{فِعْلَ النَّهْيِ:}\\
        فَفُزْ بِعِلْمٍ وَ لا تَطْلُبْ بِهِ بَدَلاً / فَالنَّاسُ مَوْتَى وَ
        أَهْلُ الْعِلْمِ أَحْيَاءُ
  \item \textbf{اِسْمَ الْفَاعِلِ:}\\
        أَيُّهَا الْفَاخِرُ جَهْلاً بِالنَّسَبِ / إِنَّمَا النَّاسُ لِأُمٍّ وَ
        لِأَبِ
  \item \textbf{اَلْفِعْلَ الْمَجْهُولَ:}\\
        هَلْ تَرَاهُم خُلِقُوا مِنْ فِضَّةٍ / أَمْ حَدِيدٍ أَمْ نُحَاسٍ أَمْ
        ذَهَبْ
  \item \textbf{اَلْفِعْلَ الْمُضَارِعَ:}\\
        بَلْ تَرَاهُم خُلِقُوا مِنْ طِينَةٍ / هَلْ سِوَى لَحْمٍ وَ عَظْمٍ وَ
        عَصَبْ
  \item \textbf{اَلْجَارَّ وَ الْمَجْرُورَ:}\\
        إِنَّمَا الْفَخْرُ لِعَقْلٍ ثَابِتٍ / وَ حَيَاءٍ وَ عَفَافٍ وَ أَدَبْ
\end{enumerate}
\end{JTamrin}

\begin{JTamrin}{٣}{جایِ خالی در ترجمه + نوعِ «لا»}
اِمْلَأِ الْفَرَاغَ فِي تَرْجَمَةِ مَا يَلِي؛ ثُمَّ عَيِّنْ نَوْعَ «لا» فِيهِ:

\begin{enumerate}
  \item ﴿وَ لا تَسُبُّوا الَّذِينَ يَدْعُونَ مِنْ دُونِ اللّٰهِ فَيَسُبُّوا
        اللّٰهَ ...﴾ اَلْأَنْعَام: ١٠٨\\
        و کسانی را که به جای خدا فرا می‌خوانند \JdotsW{4cm} که به خدا دشنام
        دهند ... .
  \item ﴿يَا أَيُّهَا الرَّسُولُ لا يَحْزُنْكَ الَّذِينَ يُسَارِعُونَ\FNmark{١}
        فِي الْكُفْرِ ...﴾ اَلْمَائِدَة: ٤١\\
        ای پیامبر، کسانی که در کفر شتاب می‌ورزند تو را \JdotsW{4cm}.
        \FNfoot{١}{سَارَعَ: شتافت}
  \item ﴿... رَبَّنَا وَ لا تُحَمِّلْنَا\FNmark{٢} مَا لا طَاقَةَ لَنَا بِهِ
        ...﴾ اَلْبَقَرَة: ٢٨٦\\
        {[ای]} پروردگارِ ما، آنچه را هیچ توانی نسبت به آن نداریم بر ما
        \JdotsW{4cm}. \FNfoot{٢}{حَمَّلَ: تحمیل کرد}
  \item لا يَرْحَمُ اللّٰهُ مَنْ لا يَرْحَمُ النَّاسَ. رَسُولُ اللّٰهِ (ص)\\
        خدا رحم نمی‌کند به کسی که به مردم \JdotsW{4cm}.
  \item يَا حَبِيبِي، لا شَيْءَ أَجْمَلُ مِنَ الْعَفْوِ عِنْدَ الْقُدْرَةِ.\\
        ای دوستِ من، \JdotsW{3cm} زیباتر از بخشش به هنگامِ توانایی نیست.
\end{enumerate}
\end{JTamrin}

\begin{JTamrin}{٤}{ترجمهٔ احادیث + تعیینِ مطلوب}
تَرْجِمِ الْأَحَادِيثَ، ثُمَّ عَيِّنِ الْمَطْلُوبَ مِنْكَ:

\begin{enumerate}
  \item كُلُّ طَعَامٍ لا يُذْكَرُ اسْمُ اللّٰهِ عَلَيْهِ ... لا بَرَكَةَ فِيهِ.
        رَسُولُ اللّٰهِ (ص)\hfill (نَائِبَ الْفَاعِلِ وَ نَوْعَ لا)
  \item لا تَجْتَمِعُ خَصْلَتَانِ فِي مُؤْمِنٍ: اَلْبُخْلُ وَ الْكَذِبُ.
        رَسُولُ اللّٰهِ (ص)\hfill (اَلْجَارَّ وَ الْمَجْرُورَ، وَ الْفَاعِلَ
        وَ إِعْرَابَهُ وَ نَوْعَ الْفِعْلِ)
  \item لا تَغْضَبْ، فَإِنَّ الْغَضَبَ مَفْسَدَةٌ\FNmark{١}. رَسُولُ اللّٰهِ (ص)
        \hfill (نَوْعَ الْفِعْلِ، وَ اسْمَ الْحَرْفِ الْمُشَبَّهِ بِالْفِعْلِ
        وَ خَبَرَهُ وَ إِعْرَابَهُمَا) \FNfoot{١}{اَلْمَفْسَدَة: مایهٔ تباهی}
  \item لا فَقْرَ أَشَدُّ مِنَ الْجَهْلِ وَ لا عِبَادَةَ مِثْلُ التَّفَكُّرِ.
        رَسُولُ اللّٰهِ (ص)\hfill (اَلْمُضَافَ إِلَيْهِ وَ نَوْعَ لا)
  \item لا تَسُبُّوا النَّاسَ فَتَكْتَسِبُوا الْعَدَاوَةَ بَيْنَهُمْ.
        رَسُولُ اللّٰهِ (ص)\hfill (نَوْعَ لا، وَ مُضَادَّ عَدَاوَة)
  \item لا تُمِيتُوا\FNmark{٢} الْقُلُوبَ بِكَثْرَةِ الطَّعَامِ وَ الشَّرَابِ
        فَإِنَّ الْقَلْبَ يَمُوتُ كَالزَّرْعِ إِذَا كَثُرَ عَلَيْهِ الْمَاءُ.
        رَسُولُ اللّٰهِ (ص)\hfill (نَوْعَ لا، وَ إِعْرَابَ الْكَلِمَاتِ
        الَّتِي تَحْتَهَا خَطٌّ: اَلْقُلُوبَ، اَلطَّعَامِ، اَلْقَلْبَ،
        اَلزَّرْعِ وَ الْمَاءُ) \FNfoot{٢}{لا تُمِيتُوا: نَمیرانید ← أَمَاتَ:
        میراند}
  \item خُذُوا\FNmark{٣} الْحَقَّ مِنْ أَهْلِ الْبَاطِلِ وَ لا تَأْخُذُوا
        الْبَاطِلَ مِنْ أَهْلِ الْحَقِّ كُونُوا\FNmark{٤} نُقَّادَ الْكَلَامِ.
        عِيسَى بْنُ مَرْيَمَ (ع)\hfill (نَوْعَ لا، وَ إِعْرَابَ الْكَلِمَاتِ
        الَّتِي تَحْتَهَا خَطٌّ: اَلْحَقَّ، اَلْبَاطِلِ، اَلْبَاطِلَ، أَهْلِ،
        اَلْكَلَامِ) \FNfoot{٣}{خُذُوا: بگیرید ← أَخَذَ: گرفت}
        \FNfoot{٤}{كُونُوا: باشید ← كَانَ: بود}
\end{enumerate}
\end{JTamrin}

\begin{JTamrin}{٥}{ترجمهٔ انواعِ فعل (بخشِ نخست)}
تَرْجِمْ أَنْوَاعَ الْفِعْلِ فِي الْجُمَلِ التَّالِيَةِ (معنای دقیقِ فعل را با
توجه به صیغه‌اش در ترجمه بیاور):

\begin{enumerate}
  \item رَجَاءً، اُكْتُبْنَ حَلَّ الْأَسْئِلَةِ مَعَ زَمِيلَاتِكُنَّ. | لطفاً
        حلِّ پرسش‌ها را با هم‌کلاسی‌هایتان \JdotsW{3.5cm}
  \item إِنْ تَكْتُبْ بِعَجَلَةٍ، فَسَيُصْبِحُ خَطُّكَ قَبِيحًا. | اگر با عجله
        \JdotsW{3.2cm} خطّت زشت خواهد شد.
  \item أُرِيدُ أَنْ أَكْتُبَ ذِكْرَى فِي دَفْتَرِ الذِّكْرَيَاتِ. | می‌خواهم
        [که] خاطره‌ای در دفترِ خاطرات \JdotsW{2.8cm}
  \item لَنْ يَكْتُبَ الْعَاقِلُ عَلَى الْآثَارِ التَّارِيخِيَّةِ. | [انسانِ]
        خردمند، رویِ آثارِ تاریخی \JdotsW{3.2cm}
  \item إِنَّهُ سَوْفَ يَكْتُبُ أَفْكَارَهُ عَلَى الْوَرَقَةِ. | قطعاً او
        اندیشه‌هایش را روی کاغذ \JdotsW{3.2cm}
  \item عِنْدَمَا شَاهَدْتُهُ كَانَ يَكْتُبُ تَمْرِينَهُ. | هنگامی‌که او را
        دیدم، تمرینش را \JdotsW{3.2cm}
\end{enumerate}
\end{JTamrin}

\begin{JTamrin}[JCompact]{٥ (ادامه)}{ترجمهٔ انواعِ فعل}
\begin{enumerate}\setcounter{enumi}{6}
  \item قَدْ يَكْتُبُ الْكَسُولُ تَمَارِينَ الدَّرْسِ. | تنبل، تمرین‌های درس را
        \JdotsW{3.8cm}
  \item يَا تَلَامِيذُ، لِمَ لا تَكْتُبُونَ التَّرْجَمَةَ؟ | ای دانش‌آموزان،
        چرا ترجمه را \JdotsW{3.2cm}
  \item رَجَاءً، لا تَكْتُبُوا عَلَى جِلْدِ الْكِتَابِ. | لطفاً رویِ جلدِ کتاب
        \JdotsW{3.2cm}
  \item هِيَ قَدْ كَتَبَتْ رِسَالَةً لِصَدِيقَتِهَا. | او نامه‌ای به دوستش
        \JdotsW{3.2cm}
  \item لِنَكْتُبْ جُمَلًا جَمِيلَةً فِي دَفَاتِرِنَا. | در دفترهایمان
        جمله‌هایی زیبا \JdotsW{3.2cm}
  \item أَنَا لَمْ أَكْتُبْ وَاجِبَاتِي أَمْسِ. | من دیروز تکلیف‌هایم را
        \JdotsW{3.2cm}
  \item مَا كَتَبْنَا شَيْئًا عَلَى الْجِدَارِ. | روی دیوار چیزی
        \JdotsW{3.2cm}
  \item كُتِبَ حَدِيثٌ عَلَى اللَّوْحِ. | حدیثی روی تخته \JdotsW{3.2cm}
  \item كَانُوا يَكْتُبُونَ الْأَجْوِبَةَ. | جواب‌ها را \JdotsW{3.2cm}
  \item كَانُوا قَدْ كَتَبُوا دَرْسَهُمْ. | درسشان را \JdotsW{3.2cm}
  \item يُكْتَبُ نَصٌّ قَصِيرٌ. | متنی کوتاه \JdotsW{3.2cm}
\end{enumerate}
\end{JTamrin}

\begin{JTamrin}{٦}{اسمِ فاعل، مفعول، مبالغه، مکان و تفضیل}
عَيِّنِ اسْمَ الْفَاعِلِ وَ اسْمَ الْمَفْعُولِ وَ اسْمَ الْمُبَالَغَةِ وَ
اسْمَ الْمَكَانِ وَ اسْمَ التَّفْضِيلِ فِي هٰذِهِ الْعِبَارَاتِ؛ ثُمَّ عَيِّنْ
تَرْجَمَةَ الْكَلِمَاتِ الْحَمْرَاءِ (در هر آیه، مشتقِّ خواسته‌شده و ترجمهٔ
واژه‌های قرمز را بیاور):

\begin{enumerate}
  \item ﴿سُبْحَانَ الَّذِي أَسْرَى\jr{بِعَبْدِهِ} لَيْلاً مِنَ الْمَسْجِدِ
        الْحَرَامِ إِلَى الْمَسْجِدِ الْأَقْصَى ...﴾ اَلْإِسْرَاء: ١
        {\JWRedWord{أَسْرَى ، لَيْلاً}}
        پاک است آن [خدایی] که بنده‌اش را شبانه از مسجدالحرام تا مسجدالاقصی
        حرکت داد.
  \item ﴿... وَ جَادِلْهُم\jr{بِالَّتِي} هِيَ أَحْسَنُ إِنَّ رَبَّكَ هُوَ
        أَعْلَمُ بِمَنْ ضَلَّ عَنْ سَبِيلِهِ ...﴾ اَلنَّحْل: ١٢٥
        {\JWRedWord{جَادِلْ ، ضَلَّ}}
        و با آنان با [شیوه‌ای] که نکوتر است بحث کن؛ قطعاً پروردگارت به [حالِ]
        کسی که از راهش گم شده، داناتر است.
  \item ﴿... يَقُولُونَ بِأَفْوَاهِهِمْ مَا لَيْسَ فِي قُلُوبِهِمْ وَ اللّٰهُ
        أَعْلَمُ بِمَا يَكْتُمُونَ\jr{...}﴾ آلِ عِمْرَان: ١٦٧
        {\JWRedWord{أَفْوَاهِ ، يَكْتُمُونَ}}
        با دهان‌هایشان (زبان‌هایشان) چیزی را می‌گویند که در دل‌هایشان نیست و
        خدا به آنچه پنهان می‌کنند داناتر است.
  \item ﴿وَ مَا أُبَرِّئُ نَفْسِي إِنَّ النَّفْسَ لَأَمَّارَةٌ\FNmark{١}
        بِالسُّوءِ إِلَّا مَا رَحِمَ رَبِّي ...﴾ يُوسُف: ٥٣
        {\JWRedWord{أُبَرِّئُ ، أَمَّارَةٌ}}
        و نفسِ خود را بی‌گناه نمی‌شمارم؛ زیرا نفس، بسیار فرمان‌ده به بدی است؛
        مگر این‌که پروردگارم رحم کند. \FNfoot{١}{آیه در کتاب با این مضمون آمده
        است.}
  \item ﴿قَدْ أَفْلَحَ الْمُؤْمِنُونَ الَّذِينَ هُمْ فِي صَلَاتِهِمْ
        خَاشِعُونَ﴾ اَلْمُؤْمِنُون: ١ و ٢ {\JWRedWord{قَدْ أَفْلَحَ ،
        خَاشِعُونَ}}
        به‌راستی‌که مؤمنان رستگار شده‌اند؛ همانان که در نمازشان فروتن‌اند.
  \item ﴿... قَالُوا لا عِلْمَ لَنَا إِنَّكَ أَنْتَ عَلَّامُ الْغُيُوبِ﴾
        اَلْمَائِدَة: ١٠٩ {\JWRedWord{عَلَّامُ}}
        گفتند: هیچ دانشی نداریم؛ قطعاً تو بسیار دانای نهان‌ها هستی.
  \item ﴿... وَ أَحْسِنُوا إِنَّ اللّٰهَ يُحِبُّ الْمُحْسِنِينَ\jr{...}﴾
        اَلْبَقَرَة: ١٩٥ {\JWRedWord{أَحْسِنُوا}}
        و نیکی کنید؛ زیرا خدا نیکوکاران را دوست می‌دارد.
  \item ﴿وَ جَعَلْنَا السَّمَاءَ سَقْفًا مَحْفُوظًا ...﴾ اَلْأَنْبِيَاء: ٣٢
        {\JWRedWord{جَعَلْنَا}}
        و آسمان را سقفی نگه‌داشته‌شده قرار دادیم.
\end{enumerate}
\end{JTamrin}

\begin{JTamrin}{٧}{جایِ خالی با حروفِ مشبههٔ فعل}
ضَعْ فِي الْفَرَاغِ كَلِمَةً مُنَاسِبَةً (گزینهٔ درست را برگزین):

\begin{enumerate}
  \item سُئِلَ الْمُدِيرُ: أَ فِي الْمَدْرَسَةِ طَالِبٌ؟ فَأَجَابَ: «
        \JdotsW{1.4cm} طَالِبَ هُنَا».\\
        (لِأَنَّ \JTick\quad لا \JTick\quad فَإِنَّ \JTick)
  \item حَضَرَ السُّيَّاحُ فِي قَاعَةِ الْمَطَارِ؛ \JdotsW{1.4cm} الدَّلِيلَ
        لَمْ يَحْضُرْ.\\
        (أَنَّ \JTick\quad لٰكِنَّ \JTick\quad لَعَلَّ \JTick)
  \item هٰذَا التَّمْرِينُ سَهْلٌ وَ \JdotsW{1.4cm} أَكْثَرُ الزُّمَلَاءِ
        خَائِفُونَ.\\
        (لَيْتَ \JTick\quad إِنَّ \JTick\quad لٰكِنَّ \JTick)
  \item قَالَ الْمُدِيرُ: \JdotsW{1.4cm} طَالِبَ رَاسِبٌ فِي الاِمْتِحَانَاتِ.\\
        (لا \JTick\quad لَنْ \JTick\quad لَمْ \JTick)
  \item تَمَنَّى الْمُزَارِعُ: «\JdotsW{1.4cm} الْمَطَرَ يَنْزِلُ كَثِيرًا!»\\
        (كَأَنَّ \JTick\quad لِأَنَّ \JTick\quad لَيْتَ \JTick)
  \item أَ لا تَعْلَمُ \JdotsW{1.4cm} الصَّبْرَ مِفْتَاحُ الْفَرَجِ؟\\
        (لٰكِنَّ \JTick\quad أَنَّ \JTick\quad لا \JTick)
  \item لِمَاذَا يَبْكِي الطِّفْلُ؟ \JdotsW{1.4cm} جَائِعٌ.\\
        (أَنَّهُ \JTick\quad لِأَنَّهُ \JTick\quad لَيْتَ \JTick)
\end{enumerate}
\end{JTamrin}

\begin{JTamrin}{٨}{جایِ خالی با اعرابِ درست}
ضَعْ فِي الْفَرَاغِ كَلِمَةً مُنَاسِبَةً (حرکتِ درستِ آخرِ واژه را برگزین):

\begin{enumerate}
  \item اِعْلَمُوا أَنَّ \JdotsW{2.2cm} النَّاسِ إِلَيْكُمْ مِنْ نِعَمِ اللّٰهِ
        عَلَيْكُمْ. اَلْإِمَامُ الْحُسَيْنُ (ع)\\
        (حَوَائِجُ \JTick\quad حَوَائِجَ \JTick\quad حَوَائِجِ \JTick)
  \item إِنَّ أَحْسَنَ الْحَسَنِ \JdotsW{2.2cm} الْحَسَنُ. اَلْإِمَامُ الْحَسَنُ
        (ع)\\
        (اَلْخُلُقُ \JTick\quad اَلْخُلُقَ \JTick\quad اَلْخُلُقِ \JTick)
  \item ﴿فَاصْبِرْ إِنَّ وَعْدَ اللّٰهِ \JdotsW{2.2cm} وَ اسْتَغْفِرْ
        لِذَنْبِكَ﴾ غَافِر: ٥٥\\
        (حَقٌّ \JTick\quad حَقًّا \JTick\quad حَقُّ \JTick)
  \item ﴿وَ أَوْفُوا بِالْعَهْدِ إِنَّ \JdotsW{2.2cm} كَانَ مَسْئُولًا﴾
        اَلْإِسْرَاء: ٣٤\\
        (اَلْعَهْدُ \JTick\quad اَلْعَهْدَ \JTick\quad اَلْعَهْدِ \JTick)
  \item ﴿إِنَّ رَحْمَتَ اللّٰهِ \JdotsW{2.2cm} مِنَ الْمُحْسِنِينَ﴾
        اَلْأَعْرَاف: ٥٦\\
        (قَرِيبٌ \JTick\quad قَرِيبًا \JTick\quad قَرِيبِ \JTick)
  \item ﴿وَ اسْتَغْفِرُوا اللّٰهَ إِنَّ اللّٰهَ \JdotsW{2.2cm} رَحِيمٌ﴾
        اَلْبَقَرَة: ١٩٩\\
        (غَفُورٌ \JTick\quad غَفُورًا \JTick\quad غَفُورِ \JTick)
  \item ﴿لا تَدْرِي لَعَلَّ \JdotsW{2.2cm} يُحْدِثُ بَعْدَ ذٰلِكَ أَمْرًا﴾
        اَلطَّلَاق: ١ \FNfoot{١}{نمی‌دانی شاید خدا پس از آن، پیشامدی نو پدید
        آورد.}
        (اَللّٰهُ \JTick\quad اَللّٰهَ \JTick\quad اَللّٰهِ \JTick)
\end{enumerate}
\end{JTamrin}

\begin{JGood}[آشنایی با ساختِ فعل‌ها: فاء، عین و لام]
در فرهنگ‌های لغتِ عربی (و در کنکور به کار می‌آید):
\begin{itemize}
  \item حرفِ اوّلِ اصلیِ فعل \textbf{فَاءُ الْفِعْل}، حرفِ دوم \textbf{عِينُ
        الْفِعْل} و حرفِ سوم \textbf{لَامُ الْفِعْل} نام دارد؛ مثال:
        يَذْهَبُ = يَفْعَلُ → يـ (حرفِ مضارع) ، ذ (فاءُ الفعل) ، هـ (عينُ
        الفعل) ، ب (لامُ الفعل).
  \item مقابلِ فعل در جدول‌های فرهنگ‌ها، حرکتِ عین‌الفعل آمده است: حرکتِ
        \textbf{فتحه} یعنی مضارعِ آن مفتوح است (ذَهَبَ ← يَذْهَبُ ؛ نَامَ ←
        يَنَامُ)؛ حرکتِ \textbf{ضمّه} یعنی مضموم است (كَتَبَ ← يَكْتُبُ ؛
        قَالَ ← يَقُولُ)؛ حرکتِ \textbf{کسره} یعنی مکسور است (جَلَسَ ←
        يَجْلِسُ ؛ سَارَ ← يَسِيرُ).
\end{itemize}
(در امتحانات و کنکور از این صفحه سؤال طرح نمی‌شود؛ اما برای شناختِ صیغه‌ها
سودمند است.)
\end{JGood}

% ============================================================================
%  ۷) پاسخ‌نامهٔ خودآزمایی و تمرین‌ها
% ============================================================================
\Jsec{اَلْإِجَابَاتُ \,\textbar\, پاسخ‌نامهٔ خودآزمایی و تمرین‌ها}

\begin{JKey}[پاسخ‌نامهٔ خودآزمایی]
\JAnswer{١}{%
  ۱. پس این روزِ رستاخیز است، ولی شما نمی‌دانستید (به آن آگاه نبودید).\\
  ۲. بی‌گمان خدا کسانی را که صف‌درصف در راهِ او جهاد می‌کنند دوست می‌دارد؛
  گویی آنان بنایی استوارند.}

\JAnswer{٢}{%
  ۱. لَيْتَ ← اسم: فَصْلَ (اسمِ لَيْتَ، منصوب به فتحه) ، خبر: طَوِيلٌ (خبرِ
  لَيْتَ، مرفوع به ضمّه) ؛ لِأَنَّ ← اسم: الرَّبِيعَ (اسمِ لِأَنَّ، منصوب) ،
  خبر: قَصِيرٌ (خبرِ لِأَنَّ، مرفوع).\\
  ۲. كَأَنَّ ← اسم: الْمُشْتَرِيَ (اسمِ كَأَنَّ، منصوب) ، خبر: مُتَرَدِّدٌ
  (مرفوع) ؛ وَلٰكِنَّ ← اسم: الْبَائِعَ (منصوب) ، خبر: عَازِمٌ (مرفوع).\\
  ۳. لَعَلَّ ← اسم: الْكَلِمَةَ (اسمِ لَعَلَّ، منصوب) ، خبر: مَكْتُوبَةٌ
  (خبرِ لَعَلَّ، مرفوع).}

\JAnswer{٣}{%
  ۱. هیچ خیری نیست در سخنی مگر با عمل. (اسْمِ لا: قَوْلٍ ؛ خَبَرُها: خَيْرَ)\\
  ۲. هیچ جهادی مانندِ جهاد با نفس نیست. (اسْمِ لا: جِهَادَ ؛ خَبَرُها حذف
  شده و جارّ و مجرور «كَجِهَادِ النَّفْسِ» بر آن عطف شده است)\\
  ۳. هیچ پوشاکی زیباتر از تندرستی نیست. (اسْمِ لا: لِبَاسَ ؛ خَبَرُها حذف
  شده است)\\
  ۴. هیچ فقری مانندِ نادانی و هیچ میراثی مانندِ ادب نیست. (اسْمِ لا: فَقْرَ
  و مِيرَاثَ ؛ خَبَرُها حذف شده است)\\
  ۵. هیچ بدی‌ای بدتر از دروغ نیست. (اسْمِ لا: سُوءَ ؛ خَبَرُها: أَسْوَأُ)}

\JAnswer{٤}{%
  ۱. نمی‌دانند؛ «لا» نفیِ مضارع است.\\
  ۲. اندوهگین نمی‌شود؛ از آنِ خدا؛ «لا»ی نخست: لای نهی به معنای «نباید»؛
  «إِنَّ»: قطعاً.\\
  ۳. هیچ تردیدی در آن نیست و؛ لا نافیهٔ جنس (اَلرَّيْبَ: اسمِ لا).\\
  ۴. هیچ اکراهی نیست؛ لا نافیهٔ جنس (إِكْرَاهَ: اسمِ لا، بدونِ ال و مفتوح).\\
  ۵. (ای مستمندان!) خوراک ندهید از آنچه خود نمی‌خورید؛ «لا» نخست: لای نهی
  (نخورید!) و «لا»ی دوم: نفیِ مضارع (نمی‌خورید).}
\end{JKey}

\begin{JKey}[پاسخ‌نامهٔ تمرین‌ها]
\JAnswer{تمرینِ ١}{%
  ۱. اَلْـ\textbf{نُّحَاس} (مس، فلزی رسانا است) ؛ ۲. اَلـ\textbf{لَّحْم}
  (گوشت) ؛ ۳. اَلـ\textbf{عَصَب} (پی/عصب) ؛ ۴. اَلـ\textbf{عَظْم} (استخوان) ؛
  ۵. اَلـ\textbf{طِّين ، طِّينَة} (گِل: خاکِ آمیخته با آب).}

\JAnswer{تمرینِ ٢}{%
  ۱. مبتدا: «دَوَاؤُكَ» — مرفوع به ضمّهٔ ظاهر.\\
  ۲. اسمِ تفضیل: «الْأَكْبَرُ» — خبرِ «إِنَّـكَ» و به‌جایِ آن، مرفوع (ضمّه)؛
  (محلِّ اعرابِ آن، خبرِ مشبهه است).\\
  ۳. خبر: «أَكْفَاءُ» — مرفوع به ضمّه.\\
  ۴. جمعِ مکسّر: «الرِّجَال» (مفرد: اَلرَّجُل) ؛ نیز «الْأَفْعَال» (مفرد:
  اَلْفِعْل) و «أَسْمَاء» (مفرد: اِسْم).\\
  ۵. فعلِ نهی: «لا تَطْلُبْ» (از تَلَبَـ/يَطْلُبُ: بطلب).\\
  ۶. اسمِ فاعل: «الْفَاخِرُ» (از فَخُرَ/يَفْخُرُ).\\
  ۷. فعلِ مجهول: «خُلِقُوا» (آفریده شدند، بابِ اِفعال، ماضی مجهول).\\
  ۸. فعلِ مضارع: «تَرَاهُم» (می‌بینی‌شان).\\
  ۹. جارّ و مجرور: «لِعَقْلٍ» (برای خردی استوار) ؛ نیز «وَ حَيَاءٍ وَ
  عَفَافٍ وَ أَدَبْ» بر آن عطف شده‌اند.}

\JAnswer{تمرینِ ٣}{%
  ۱. دشنام ندهید؛ «لا» نهی (بر سر فعلِ مضارعِ مخاطب، علامتِ سکون).\\
  ۲. اندوهگین نسازی؛ «لا» نفیِ مضارع (پیش از ضمیرِ مخاطب، فتحه).\\
  ۳. تحمیل نکن؛ «لا» به معنای «نباید» (بر سر فعلِ مضارع، علامتِ سکون)؛ در
  «مَا لا طَاقَةَ لَنَا بِهِ»: لا نافیهٔ جنس است (طَاقَةَ: اسمِ لا، بدونِ ال
  و مفتوح).\\
  ۴. رحم نمی‌کند؛ «لا»ی نخست: نفیِ مضارع؛ «لا»ی دوم: نفیِ مضارع.\\
  ۵. هیچ چیزی ... نیست؛ «لا» نافیهٔ جنس (شَيْءَ: اسمِ لا).}

\JAnswer{تمرینِ ٤}{%
  ۱. هر خوراکی که نامِ خدا بر آن گفته نشود، در آن هیچ برکتی نیست.
  (نائبِ فاعل: «اسْمُ اللّٰهِ» ؛ نوعِ لا: نافیهٔ جنس)\\
  ۲. دو خوی در مؤمن گرد نمی‌آید: بخل و دروغ. (جارّ و مجرور: «فِي مُؤْمِنٍ» ؛
  فاعل: «خَصْلَتَانِ»، مرفوع به ضمّهٔ ظاهر؛ فعل: ماضی، معلوم)\\
  ۳. خشمگین مباش که خشم، مایهٔ تباهی است. (فعل: نهی «لا تَغْضَبْ» ؛ اسمِ
  مشبهه: «الْغَضَبَ» منصوب ؛ خبر: «مَفْسَدَةٌ» مرفوع)\\
  ۴. هیچ فقری سخت‌تر از نادانی و هیچ عبادتی مانندِ تفکر نیست. (مضافٌ‌الیه:
  «الْجَهْلِ» و «التَّفَكُّرِ» مجرور ؛ نوعِ هر دو «لا»: نافیهٔ جنس)\\
  ۵. مردم را دشنام ندهید که دشمنی میانشان می‌اندوزید. (نوعِ لا: نهی ؛
  متضادِّ عداوة: اَلْمَوَدَّة/ اَلْمَحَبَّة)\\
  ۶. دل‌ها را با فراوانیِ خوراک و نوشیدنی نمیرانید که قلب می‌میرد مانندِ
  کشتزار وقتی آبِ فراوان بر آن ببارد. (نوعِ لا: نهی ؛ اعراب: اَلْقُلُوبَ
  مفعول، فتحه؛ اَلطَّعَامِ مضافٌ‌الیه، کسره؛ اَلْقَلْبَ مبتدا یا اسمِ إنّ،
  فتحه؛ اَلزَّرْعِ مجرور بعد از «كَـ»، کسره؛ اَلْمَاءُ فاعل، مرفوع به ضمّه)\\
  ۷. حق را از اهلِ باطل بگیرید و باطل را از اهلِ حق نگیرید؛ نقدکنندهٔ
  سخنان باشید. (نوعِ لا: نهی ؛ اَلْحَقَّ مفعول، فتحه؛ اَلْبَاطِلِ مجرور بعد
  از «مِنْ»، کسره؛ اَلْبَاطِلَ مفعول، فتحه؛ أَهْلِ مضافٌ‌الیه، کسره؛
  اَلْكَلَامِ مضافٌ‌الیه، کسره)}

\JAnswer{تمرینِ ٥}{%
  ۱. بنویسید (امرِ مؤنثِ جمع) ؛ ۲. بنویسی (مضارعِ مجزومِ بعد از إِنْ
  شرطیّه) ؛ ۳. بنویسم (مضارعِ منصوبِ بعد از أنْ) ؛ ۴. نخواهد نوشت (لَنْ +
  مضارعِ منصوب) ؛ ۵. خواهد نوشت (سَـ/سَوْفَ + مضارع) ؛ ۶. می‌نوشت (ماضیِ
  استمراری: كَانَ + مضارع) ؛ ۷. شاید بنویسد (قَدْ + مضارع = احتمال) ؛
  ۸. نمی‌نویسید (لا + مضارع = نفیِ مضارع) ؛ ۹. ننویسید (لا + مضارع = نهی) ؛
  ۱۰. نوشته است (قَدْ + ماضی = گذشتهٔ نقلی) ؛ ۱۱. بنویسیم (لِـ + مضارعِ
  مجزوم = امر به لام) ؛ ۱۲. ننوشتم (لَمْ + مضارعِ مجزوم = ماضیِ منفی) ؛
  ۱۳. ننوشتیم (مَا + ماضی = نفیِ ماضی) ؛ ۱۴. نوشته شد (ماضیِ مجهول) ؛
  ۱۵. می‌نوشتند (كَانُوا + مضارع = ماضیِ استمراری) ؛ ۱۶. نوشته بودند (كَانُوا
  قَدْ + ماضی = ماضیِ بعید/التزامی) ؛ ۱۷. نوشته می‌شود (مضارعِ مجهول).}

\JAnswer{تمرینِ ٦}{%
  ۱. اسمِ فاعل: أَسْرَى (از سَرَی) ؛ ترجمهٔ قرمز: بنده‌اش را ، شبانه.\\
  ۲. امرِ مخاطبِ مفرد: جَادِلْ ؛ اسمِ تفضیل: أَحْسَنُ و أَعْلَمُ ؛ ترجمهٔ
  قرمز: بحث کن ؛ گم شده، راهِ خود را باخته.\\
  ۳. جمعِ مکسّر: أَفْوَاه (مفرد: فَم) ؛ فعلِ مضارع: يَكْتُمُونَ ؛ ترجمهٔ
  قرمز: دهان‌ها ؛ پنهان می‌کنند.\\
  ۴. اسمِ مبالغه: أَمَّارَةٌ (صیغهٔ فَعَّالَة) ؛ فعل مضارع: أُبَرِّئُ ؛ ترجمهٔ
  قرمز: بی‌گناه نمی‌شمارم ؛ بسیار فرمان‌ده به بدی.\\
  ۵. ماضی + قَدْ: قَدْ أَفْلَحَ (رستگار شده‌اند) ؛ اسمِ فاعلِ جمع: خَاشِعُونَ
  (جمعِ مذکرِ سالم) ؛ ترجمهٔ قرمز: رستگار شده‌اند ؛ فروتن.\\
  ۶. اسمِ مبالغه: عَلَّامُ (صیغهٔ فَعَّال) ؛ ترجمهٔ قرمز: بسیار دانا.\\
  ۷. امر: أَحْسِنُوا ؛ اسمِ فاعل: الْمُحْسِنِينَ ؛ ترجمهٔ قرمز: نیکی کنید ؛
  نیکوکاران.\\
  ۸. ماضی: جَعَلْنَا (قرار دادیم) ؛ اسمِ مفعول: مَحْفُوظًا ؛ ترجمهٔ قرمز:
  قرار دادیم ؛ نگه‌داشته‌شده.}

\JAnswer{تمرینِ ٧}{%
  ۱. لا (لا نافیهٔ جنس: هیچ دانش‌آموزی اینجا نیست) ؛ ۲. لٰكِنَّ (ولی راهنما
  نیامده بود) ؛ ۳. لٰكِنَّ (این تمرین آسان است، ولی بیشترِ هم‌کلاسی‌ها
  ترسان‌اند) ؛ ۴. لا (لا نافیهٔ جنس: هیچ دانش‌آموزی مردود نیست) ؛ ۵. لَيْتَ
  (کاش بارانِ فراوان ببارد!) ؛ ۶. أَنَّ (آیا نمی‌دانی \emph{که} صبر کلید
  گشایش است؟) ؛ ۷. لِأَنَّهُ (زیرا گرسنه است).}

\JAnswer{تمرینِ ٨}{%
  ۱. حَوَائِجَ (اسمِ أَنَّ، منصوب) ؛ ۲. اَلْخُلُقُ (خبرِ إِنَّ، مرفوع) ؛
  ۳. حَقٌّ (خبرِ إِنَّ، مرفوع) ؛ ۴. اَلْعَهْدُ (مبتدا و خبرِ «كَانَ»
  مرفوع؛ نیز اسمِ إنّ) ؛ ۵. قَرِيبٌ (خبرِ إِنَّ، مرفوع) ؛ ۶. غَفُورٌ (خبرِ
  إِنَّ، مرفوع) ؛ ۷. اَللّٰهُ (اسمِ لَعَلَّ، مرفوع؛ زیرا فعلِ بعدی، فاعل
  دارد و اسمِ مشبهه با «ال» به کار می‌رود).}
\end{JKey}
